\chapter{Continuous retention diagnostic}\label{app:continuous}
The four-target diagnostic changes a real-valued coefficient on one target branch. From current real parameters and Adam moments it creates a virtual update using two older complete arrived days $A$. A newer complete arrived day $V$, not yet used in real training, supplies the forecast validation loss. Following the one-step hypergradient approach \cite{luketina2016}, the diagnostic computes
\begin{equation}\begin{aligned}
\widetilde\theta_t(\lambda)&=\operatorname{AdamStep}
\left(\theta_t,m_t,v_t;\nabla_\theta[F_A+\lambda R_A]\right),\\
g_t&=\left.\frac{\partial F_V(\widetilde\theta_t(\lambda))}
{\partial\lambda}\right|_{\lambda_t},\\
\lambda_{t+1}&=\Pi_{[0,U]}\left(\lambda_t-\eta_\lambda g_t\right).
\end{aligned}\label{eq:hypergradient}\end{equation}
Previous parameters and moments are fixed starting state for this local derivative. The differentiable virtual Adam step uses bias-corrected moments and does not mutate the real optimizer. After the decision, real training consumes the newly arrived example; the next forecast uses only permitted data.

The fixed search defines $U=33.84425$ and supplies its minimizer as $\lambda_0=0.798445$. Calibration compares three nominal step fractions at two decision cadences and selects $\eta_\lambda=211.295$ with an update at each eligible arrival. Across 801 forecast records for the continuous branch, the deployed coefficients ranged from 0 to 7.73128. The coefficient step size differs from the model learning rate 0.001. Adam uses betas $(0.9,0.999)$ and optimizer epsilon $10^{-8}$, separate from retention epsilon 0.001. The controller upper bound is the smallest measured coefficient above the selected Fixed value that is harmful relative to zero; if none is measured, it uses the largest finite tested coefficient. Here the fallback gives the stated upper endpoint. Step sizes are scaled from a calibration-prefix hypergradient pilot and selected across three nominal step fractions and two cadences. The projection interval is a study-specific bounded domain, not a paper-established EV constant or a measured claim that every value in it is useful.

